\subsection{\texorpdfstring{ALGEBRAS OF TYPE \(\mathbf{B}_l\) ( \(l \geq 1\) )}{ALGEBRAS OF TYPE \textbackslash mathbf\{B\}\_l ( l \textbackslash geq 1 )}}

\begin{enumerate}
\def\labelenumi{(\Roman{enumi})}
\tightlist
\item
  Let V be a finite dimensional vector space, and \(\Psi\) a non-degenerate symmetric bilinear form on V. The set of endomorphisms x of V such that \(\Psi(xv,v') + \Psi(v, xv') = 0\) for all \(v, v' \in V\) is a Lie subalgebra of \(\mathfrak{sl}(V)\) , semi-simple for dim \(V \neq 2\) (Chap. I, §6, no. 7, Prop. 9). We denote it by \(\mathfrak{o}(\Psi)\) and call it the orthogonal Lie algebra associated to \(\Psi\) .
\end{enumerate}

Assume that V is of odd dimension \(2l + 1 \geq 3\) and that \(\Psi\) is of maximum index l. Denote by Q the quadratic form such that \(\Psi\) is associated to Q. We have \(\mathrm{Q}(x)=\frac{1}{2}\varPsi(x,x)\) for \(x\in V\) . By Algebra, Chap. IX, §4, no. 2, V can be written as the direct sum of two maximal totally isotropic subspaces F and \(F'\) and the orthogonal complement G of \(F+F'\) , which is non-isotropic and 1-dimensional. Up to multiplying \(\varPsi\) by a non-zero constant, we can assume that there exists \(e_{0}\in G\) such that \(\varPsi(e_{0},e_{0})=-2\) . On the other hand, F and \(F'\) are in duality via \(\varPsi\) ; let \((e_{i})_{1\leq i\leq l}\) be a basis of F and \((e_{-i})_{1\leq i\leq l}\) the dual basis of \(F'\) . Then

\[
(e _ {1}, \ldots , e _ {l}, e _ {0}, e _ {- l}, \ldots , e _ {- 1})
\]

is a basis of V; we have

\[
\mathrm{Q} \left(\sum x _ {i} e _ {i}\right) = - x _ {0} ^ {2} + \sum_ {i = 1} ^ {i = l} x _ {i} x _ {- i}
\]

and the matrix of \(\varPsi\) with respect to this basis is the square matrix of order \(2l + 1\)

\[
S = \left( \begin{array}{c c c} 0 & 0 & s \\ 0 & - 2 & 0 \\ s & 0 & 0 \end{array} \right), \qquad s = \left( \begin{array}{c c c c c} 0 & 0 & \dots & 0 & 1 \\ 0 & 0 & \dots & 1 & 0 \\ \vdots & \vdots & \ddots & \vdots & \vdots \\ 0 & 1 & \dots & 0 & 0 \\ 1 & 0 & \dots & 0 & 0 \end{array} \right),
\]

where s is the square matrix of order l all of whose entries are zero except those on the second diagonal \(^{6}\) which are equal to 1. A basis of V with the preceding properties will be called a Witt basis of V. The algebra \(\mathfrak{g} = \mathfrak{o}(\varPsi)\) can then be identified with the algebra \(\mathfrak{o}_{S}(2l + 1, k)\) of square matrices a of order \(2l + 1\) such that \(a = -S^{-1}aS\) (Algebra, Chap. IX, §1, no. 10, formulas (50)). An easy calculation shows that g is the set of matrices of the form

\[
\left( \begin{array}{c c c} A & 2 s ^ {t} x & B \\ y & 0 & x \\ C & 2 s ^ {t} y & D \end{array} \right)\tag{2}
\]

where x and y are matrices with 1 row and l columns and A, B, C, D are square matrices of order l such that \(B = -s^{t}Bs\) , \(C = -s^{t}Cs\) , and \(D = -s^{t}As\) . Since the map \(A \mapsto s^{t}As\) from \(\mathbf{M}_{l}(k)\) to itself is the symmetry with respect to the second diagonal, it follows that

\[
\dim \mathfrak {g} = 2 l + l ^ {2} + 2 \frac {l (l - 1)}{2} = l (2 l + 1).
\]

Let \(\mathfrak{h}\) be the set of diagonal elements of \(\mathfrak{g}\) . This is a commutative subalgebra of \(\mathfrak{g}\) , with basis the elements

\footnotetext[6]{The second diagonal of a square matrix \((a_{ij})_{1\leq i,j\leq n}\) is the family of \(a_{ij}\) such that \(i+j=n+1\) .}

\[
H _ {i} = E _ {i, i} - E _ {- i, - i} \quad (1 \leq i \leq l).
\]

Let \((\varepsilon_{i})\) be the basis of \(h^{*}\) dual to \((H_{i})\) . Put

\[
\left\{ \begin{array}{l l l} X _ {\varepsilon_ {i}} & = 2 E _ {i, 0} + E _ {0, - i} & (1 \leq i \leq l) \\ X _ {- \varepsilon_ {i}} & = - 2 E _ {- i, 0} - E _ {0, i} & (1 \leq i \leq l) \\ X _ {\varepsilon_ {i} - \varepsilon_ {j}} & = E _ {i, j} - E _ {- j, - i} & (1 \leq i <   j \leq l) \\ X _ {\varepsilon_ {j} - \varepsilon_ {i}} & = - E _ {j, i} + E _ {- i, - j} & (1 \leq i <   j \leq l) \\ X _ {\varepsilon_ {i} + \varepsilon_ {j}} & = E _ {i, - j} - E _ {j, - i} & (1 \leq i <   j \leq l) \\ X _ {- \varepsilon_ {i} - \varepsilon_ {j}} & = - E _ {- j, i} + E _ {- i, j} & (1 \leq i <   j \leq l). \end{array} \right.\tag{3}
\]

It is easy to verify that these elements form a basis of a complement of \(\mathfrak{h}\) in \(\mathfrak{g}\) and that, for \(h\in \mathfrak{h}\) ,

\[
[ h, X _ {\alpha} ] = \alpha (h) X _ {\alpha}\tag{4}
\]

for all \(\alpha \in \mathbb{R}\) , where \(\mathbb{R}\) is the set of the \(\pm \varepsilon_i\) and the \(\pm \varepsilon_i \pm \varepsilon_j\) ( \(1 \leq i < j \leq l\) ). It follows that \(\mathfrak{h}\) is equal to its normalizer in \(\mathfrak{g}\) , and hence is a Cartan subalgebra of \(\mathfrak{g}\) , that \(\mathfrak{h}\) is splitting, and that the roots of \((\mathfrak{g}, \mathfrak{h})\) are the elements of \(\mathbb{R}\) . The root system \(\mathbb{R}\) of \((\mathfrak{g}, \mathfrak{h})\) is of type \(\mathrm{B}_l\) for \(l \geq 2\) , and of type \(\mathrm{A}_1\) (also said to be of type \(\mathrm{B}_1\) ) for \(l = 1\) (Chap. VI, §4, no. 5.I, extended to the case \(l = 1\) ). Consequently, \(\mathfrak{g}\) is a splittable simple Lie algebra of type \(\mathrm{B}_l\) .

Every splitting Cartan subalgebra of \(\mathfrak{o}(\varPsi)\) is a transform of h by an elementary automorphism of \(\mathfrak{o}(\varPsi)\) , and hence by an element of \(\mathbf{O}(\varPsi)\) (cf.~(VII)), and consequently is the set \(h_{\beta}\) of elements of g whose matrix with respect to a Witt basis \(\beta\) of V is diagonal. We verify immediately that the only vector subspaces invariant under \(h_{\beta}\) are those generated by a subset of \(\beta\) .

If \(l = 1\) , the algebras \(\mathfrak{o}(\varPsi)\) and \(\mathfrak{sl}(2, k)\) have the same root systems, and are thus isomorphic (cf.~also §1, Exerc. 16). From now on, we assume that \(l \geq 2\) .

\begin{enumerate}
\def\labelenumi{(\Roman{enumi})}
\setcounter{enumi}{1}
\tightlist
\item
  The root system \(R^{\vee}\) is determined by means of Chap. VI, §4, no. 5.V, and we find that
\end{enumerate}

\[
H _ {\varepsilon_ {i}} = 2 H _ {i}, \quad H _ {\varepsilon_ {i} - \varepsilon_ {j}} = H _ {i} - H _ {j}, \quad H _ {\varepsilon_ {i} + \varepsilon_ {j}} = H _ {i} + H _ {j}.
\]

\begin{enumerate}
\def\labelenumi{(\Roman{enumi})}
\setcounter{enumi}{2}
\tightlist
\item
  Put \(\alpha_{1} = \varepsilon_{1} - \varepsilon_{2},\ldots ,\alpha_{l - 1} = \varepsilon_{l - 1} - \varepsilon_{l},\alpha_{l} = \varepsilon_{l}\) . By Chap. VI, §4, no. 5.II, \((\alpha_{1},\dots ,\alpha_{l})\) is a basis B of R; the positive roots relative to B are the \(\varepsilon_{i}\) and the \(\varepsilon_{i}\pm \varepsilon_{j}\) ( \(i < j\) ). The corresponding Borel subalgebra \(\mathfrak{b}\) is the set of upper triangular matrices in \(\mathfrak{g}\) .
\end{enumerate}

It is immediately verified that the only vector subspaces of V distinct from \(\{0\}\) and V stable under b are the elements of the maximal flag corresponding to the basis \((e_{i})\) , that is, the totally isotropic subspaces \(V_{1},\ldots,V_{l}\) , where \(V_{i}\) is generated by \(e_{1},\ldots,e_{i}\) , together with their orthogonal complements \(V_{-1},\ldots,V_{-i}\) : the orthogonal complement \(V_{-i}\) of \(V_{i}\) is generated by \(e_{1},\ldots,e_{l},e_{0},e_{-l},\ldots,e_{-i-1}\) and is not totally isotropic. On the other hand, if an element of g leaves stable a vector subspace, it leaves stable its orthogonal complement. Consequently, \(\mathfrak{b}\) is the set of elements of \(\mathfrak{g}\) leaving stable the elements of the flag \(\{\mathrm{V}_1,\dots ,\mathrm{V}_l\}\) .

A flag is said to be isotropic if each of its elements is totally isotropic. The flag \(\{V_{1},\ldots,V_{l}\}\) is a maximal isotropic flag. Since the group \(\mathbf{O}(\varPsi)\) operates transitively both on the Borel subalgebras of g (cf.~(VII)) and on the maximal isotropic flags (Algebra, Chap. IX, §4, no. 3, Th. 1), we see that, for any maximal isotropic flag \(\delta\) in V, the set \(b_{\delta}\) of elements of g leaving stable the elements of \(\delta\) is a Borel subalgebra of g and that the map \(\delta\mapsto b_{\delta}\) is a bijection from the set of maximal isotropic flags to the set of Borel subalgebras of g.

Let \(\delta\) be an isotropic flag and let \(\mathfrak{p}_{\delta}\) be the set of elements of \(\mathfrak{g}\) leaving stable the elements of \(\delta\) . If \(\delta \subset \{\mathrm{V}_1, \ldots, \mathrm{V}_l\}\) , then \(\mathfrak{p}_{\delta}\) is a parabolic subalgebra of \(\mathfrak{g}\) containing \(\mathfrak{b}\) , and it is easy to verify that the only totally isotropic subspaces \(\neq \{0\}\) stable under \(\mathfrak{p}_{\delta}\) are the elements of \(\delta\) . This gives \(2^l\) parabolic subalgebras of \(\mathfrak{g}\) containing \(\mathfrak{b}\) . We see as above that the map \(\delta \mapsto \mathfrak{p}_{\delta}\) is a bijection from the set of isotropic flags in V to the set of parabolic subalgebras of \(\mathfrak{g}\) . Moreover, \(\mathfrak{p}_{\delta} \subset \mathfrak{p}_{\delta'}\) if and only if \(\delta \supset \delta'\) .

\begin{enumerate}
\def\labelenumi{(\Roman{enumi})}
\setcounter{enumi}{3}
\tightlist
\item
  The fundamental weights corresponding to \(\alpha_{1},\ldots,\alpha_{l}\) are, by Chap. VI, §4, no. 5.VI,
\end{enumerate}

\[
\begin{array}{l} \varpi_ {i} = \varepsilon_ {1} + \dots + \varepsilon_ {i} (1 \leq i \leq l - 1) \\ \varpi_ {l} = \frac {1}{2} (\varepsilon_ {1} + \dots + \varepsilon_ {l}). \end{array}
\]

Let \(\sigma\) be the identity representation of \(\mathfrak{g}\) on V. The exterior power \(\bigwedge^{r}\sigma\) operates on \(E = \bigwedge^{r}V\) . If \(h\in \mathfrak{h}\) ,

\[
\begin{array}{l} \sigma (h). e _ {i} = \varepsilon_ {i} (h) e _ {i} \quad \text { for } 1 \leq i \leq l \\ \sigma (h). e _ {0} = 0 \\ \sigma (h). e _ {- i} = - \varepsilon_ {i} (h) e _ {- i} \quad \text { for } 1 \leq i \leq l. \end{array}
\]

It follows that, for \(1 \leq r \leq l\) , \(\varepsilon_{1} + \cdots + \varepsilon_{r}\) is the highest weight of \(\bigwedge^{r}\sigma\) , the elements of weight \(\varepsilon_{1} + \cdots + \varepsilon_{r}\) being those proportional to \(e_{1} \wedge \cdots \wedge e_{r}\) . We shall show that for \(1 \leq r \leq l - 1\) , the representation \(\bigwedge^{r}\sigma\) is a fundamental representation of g of highest weight \(\varpi_{r}\) . For this, it is enough to show that \(\bigwedge^{r}\sigma\) is irreducible for \(0 \leq r \leq 2l + 1\) . But the bilinear form \(\Phi\) on \(\bigwedge^{r}V \times \bigwedge^{2l+1-r}V\) defined by

\[
x \wedge y = \varPhi (x, y) e _ {1} \wedge \dots \wedge e _ {l} \wedge e _ {0} \wedge e _ {- l} \wedge \dots \wedge e _ {- 1}
\]

is invariant under g and puts \(\bigwedge^{r}V\) and \(\bigwedge^{2l+1-r}V\) in duality. Thus, the representation \(\bigwedge^{2l+1-r}\sigma\) is the dual of \(\bigwedge^{r}\sigma\) and it suffices to prove the irreducibility of \(\bigwedge^{r}\sigma\) for \(0\leq r\leq l\) , or that the smallest subspace \(T_{r}\) of \(\bigwedge^{r}V\) containing \(e_{1}\wedge\cdots\wedge e_{r}\) and stable under g is the whole of \(\bigwedge^{r}V\) . This is immediate for r=0 and r=1 (cf.~formula (2)). For r=2 (and hence \(l\geq2\) ), the representation \(\bigwedge^{2}\sigma\) and the adjoint representation of g (which is irreducible) have the same dimension \(l(2l + 1)\) and the same highest weight \(\varepsilon_1 + \varepsilon_2\) (Chap. VI, §4, no. 5.IV). We conclude that \(\bigwedge^2\sigma\) is equivalent to the adjoint representation, and hence is irreducible. This proves our assertion for \(l = 1\) and \(l = 2\) .

We now argue by induction on \(l\) , and assume that \(l \geq r \geq 3\) . We remark first of all that if W is a non-isotropic subspace of V of odd dimension, with orthogonal complement \(\mathrm{W}'\) , the restriction \(\Psi_{\mathrm{W}}\) of \(\Psi\) to W is non-degenerate and \(\mathfrak{o}(\Psi_{\mathrm{W}})\) can be identified with the subalgebra of \(\mathfrak{g}\) consisting of the elements vanishing on \(\mathrm{W}'\) . If \(\dim \mathrm{W} < \dim \mathrm{V}\) , and if \(\Psi_{\mathrm{W}}\) is of maximal index, the induction hypothesis implies that if \(\mathrm{T}_r\) contains a non-zero element of the form \(w' \wedge w\) , with \(w' \in \bigwedge^{r - k} \mathrm{W}'\) and \(w \in \bigwedge^k \mathrm{W}\) ( \(0 \leq k \leq r\) ), then \(\mathrm{T}_r\) contains \(w' \wedge \bigwedge^k \mathrm{W}\) : indeed, we have \(a.(w' \wedge w) = w' \wedge a.w\) for all \(a \in \mathfrak{o}(\Psi_{\mathrm{W}})\) . We show by induction on \(p \in [0, r]\) that \(\mathrm{T}_r\) contains the elements

\[
x = e _ {i _ {1}} \wedge \dots \wedge e _ {i _ {r - p}} \wedge e _ {j _ {1}} \wedge \dots \wedge e _ {j _ {p}}
\]

for \(1 \leq i_{1} < \cdots < i_{r-p} \leq l\) and \(-l \leq j_{1} < \cdots < j_{p} \leq 0\) . For p = 0, this follows from the irreducibility of the operation of \(\mathfrak{gl}(F)\) on \(\bigwedge^{r} F\) (no. 1), since g contains the elements leaving \(F = V_{l} = \sum_{i=1}^{l} ke_{i}\) fixed and inducing on it any endomorphism (cf.~formula (2)). If p = 1, let \(q \in (1, l)\) be such that \(q \neq -j_{1}\) and such that there exists \(\lambda \in [1, r - p]\) with \(q = i_{\lambda}\) ; if \(p \geq 2\) , let \(q \in [1, l]\) be such that \(-q \in \{j_{1}, \ldots, j_{p}\}\) . Permuting the \(e_{i}\) if necessary, we can assume that q = 1. Now take for W the orthogonal complement of \(W' = ke_{1} + ke_{-1}\) . If p = 1, we have \(x \in e_{1} \wedge \bigwedge^{r-1} W\) ; since \(T_{r}\) contains \(e_{1} \wedge \cdots \wedge e_{r}\) , we see that \(T_{r}\) contains x. If \(p \geq 2\) , either \(x \in e_{-1} \wedge \bigwedge^{r-1} W\) or \(x \in e_{1} \wedge e_{-1} \wedge \bigwedge^{r-2} W\) ; since \(T_{r}\) contains \(e_{-1} \wedge e_{2} \wedge \cdots \wedge e_{r-1}\) and \(e_{-1} \wedge e_{1} \wedge e_{2} \wedge \cdots \wedge e_{r-2}\) by the induction hypothesis, we see that \(T_{r}\) contains x, which completes the proof.

For another proof of the irreducibility of \(\bigwedge^{r}\sigma\) , see Exerc. 6.

We shall now determine the fundamental representation with highest weight \(\varpi_{l}\) .

Lemma 1. Let \(\mathrm{V}\) be a finite dimensional vector space, \(\mathbb{Q}\) a non-degenerate quadratic form on \(\mathrm{V}\) , \(\Psi\) the symmetric bilinear form associated to \(\mathbb{Q}\) , \(\mathrm{C(Q)}\) the Clifford algebra of \(\mathrm{V}\) relative to \(\mathbb{Q}\) , \(f_0\) the composite of the canonical maps

\[
\mathfrak {o} (\Psi) \longrightarrow \mathfrak {g l} (V) \longrightarrow V \otimes V ^ {*} \longrightarrow V \otimes V \longrightarrow C ^ {+} (Q)
\]

(the 1st is the canonical injection, the 3rd is defined by the canonical isomorphism from \(\mathrm{V}^*\) to \(\mathrm{V}\) corresponding to \(\Psi\) , the 4th is defined by the multiplication in \(\mathrm{C(Q)}\) , cf.~Algebra, Chap. IX, § 9, no. 1). Put \(f = \frac{1}{2} f_0\) .

\begin{enumerate}
\def\labelenumi{(\roman{enumi})}
\item
  If \((e_r), (e_r')\) are bases of V such that \(\Psi(e_r, e_s') = \delta_{rs}\) , we have \(f_0(a) = \sum_r (ae_r)e_r'\) for all \(a \in \mathfrak{o}(\Psi)\) .
\item
  If \(a, b \in \mathfrak{o}(\Psi)\) , we have \(\sum_{r}(ae_{r})(be_{r}^{\prime}) = -\sum_{r}(abe_{r})e_{r}^{\prime}\) .
\item
  If \(a \in \mathfrak{o}(\varPsi)\) and \(v \in V\) , we have \([f(a), v] = av\) .
\item
  If \(a, b \in \mathfrak{o}(\varPsi)\) , we have \([f(a), f(b)] = f([a, b])\) .
\item
  \(f(\mathfrak{o}(\Psi))\) generates the associative algebra \(\mathrm{C}^{+}(\mathrm{Q})\) .
\item
  Let \(\mathrm{N}\) be a left \(\mathrm{C}^{+}(\mathrm{Q})\) -module and \(\rho\) the corresponding homomorphism from \(\mathrm{C}^{+}(\mathrm{Q})\) to \(\operatorname{End}_k(\mathrm{N})\) . Then \(\rho \circ f\) is a representation of \(\mathfrak{o}(\Psi)\) on \(\mathrm{N}\) . If \(\mathrm{N}\) is simple, \(\rho \circ f\) is irreducible.
\end{enumerate}

Assertion (i) is clear. If \(a, b \in \mathfrak{o}(\varPsi)\) , we have (putting \(\varPsi(x,y) = \langle x,y\rangle\) ):

\[
\begin{array}{c} \sum_ {r} (a e _ {r}) (b e _ {r} ^ {\prime}) = \sum_ {r, s, t} \langle a e _ {r}, e _ {s} ^ {\prime} \rangle \langle b e _ {r} ^ {\prime}, e _ {t} \rangle e _ {s} e _ {t} ^ {\prime} = \sum_ {r, s, t} \langle e _ {r}, a e _ {s} ^ {\prime} \rangle \langle e _ {r} ^ {\prime}, b e _ {t} \rangle e _ {s} e _ {t} ^ {\prime} \\ = \sum_ {s, t} \langle a e _ {s} ^ {\prime}, b e _ {t} \rangle e _ {s} e _ {t} ^ {\prime} = - \sum_ {s, t} \langle e _ {s} ^ {\prime}, a b e _ {t} \rangle e _ {s} e _ {t} ^ {\prime} = - \sum_ {t} (a b e _ {t}) e _ {t} ^ {\prime} \end{array}
\]

which proves (ii). Next, for all \(v \in V\) , we have by (i),

\[
\begin{array}{l} [ f (a), v ] = \frac {1}{2} \sum_ {r} ((a e _ {r}) e _ {r} ^ {\prime} v - v (a e _ {r}) e _ {r} ^ {\prime}) \\ \qquad = \frac {1}{2} \sum_ {r} ((a e _ {r}) e _ {r} ^ {\prime} v + (a e _ {r}) v e _ {r} ^ {\prime} - (a e _ {r}) v e _ {r} ^ {\prime} - v (a e _ {r}) e _ {r} ^ {\prime}) \\ \qquad = \frac {1}{2} \sum_ {r} ((a e _ {r}) \langle e _ {r} ^ {\prime}, v \rangle - \langle a e _ {r}, v \rangle e _ {r} ^ {\prime}) \\ \qquad = \frac {1}{2} a \left(\sum_ {r} \langle e _ {r} ^ {\prime}, v \rangle e _ {r}\right) + \frac {1}{2} \sum_ {r} \langle e _ {r}, a v \rangle e _ {r} ^ {\prime} = \frac {1}{2} a v + \frac {1}{2} a v = a v, \end{array}
\]

which proves (iii). Then

\[
\begin{array}{l} [ f (a), f (b) ] = \left[ f (a), \frac {1}{2} \sum_ {r} (b e _ {r}) e _ {r} ^ {\prime} \right] \quad \text { by   (i) } \\ \quad = \frac {1}{2} \sum_ {r} ([ f (a), b e _ {r} ] e _ {r} ^ {\prime} + (b e _ {r}) [ f (a), e _ {r} ^ {\prime} ]) \\ \quad = \frac {1}{2} \sum_ {r} ((a b e _ {r}) e _ {r} ^ {\prime} + (b e _ {r}) (a e _ {r} ^ {\prime})) \quad \text { by   (iii) } \\ \quad = \frac {1}{2} \sum_ {r} ((a b e _ {r}) e _ {r} ^ {\prime} - (b a e _ {r}) e _ {r} ^ {\prime}) \quad \text { by   (ii) } \\ \quad = f ([ a, b ]) \quad \text { by   (i) } \end{array}
\]

which proves (iv). To prove (v), we can, by extending scalars, assume that \(k\) is algebraically closed. Choose then a basis \((e_r)\) of V such that \(\varPsi(e_r, e_s) = \delta_{rs}\) , so \(e_r' = e_r\) . If \(i \neq j\) , then \(\mathrm{E}_{ij} - \mathrm{E}_{ji} \in \mathfrak{o}(\varPsi)\) and

\[
f (\mathrm{E} _ {i j} - \mathrm{E} _ {j i}) = \frac {1}{2} (e _ {i} e _ {j} - e _ {j} e _ {i}) = e _ {i} e _ {j};
\]

but the \(e_{i}e_{j}\) generate \(\mathrm{C}^{+}(\mathrm{Q})\) .

Assertion (vi) follows from (iv) and (v).

Q.E.D.

Recall now the notations used at the beginning of this number. Put \(\tilde{V} = F + F'\) and let \(\tilde{Q}\) (resp. \(\tilde{\Psi}\) ) be the restriction of Q (resp. \(\Psi\) ) to \(\tilde{V}\) . Then \(\tilde{Q}\) is a non-degenerate quadratic form of maximum index l on the space \(\tilde{V}\) of dimension 2l and the Clifford algebra \(C(\tilde{Q})\) is a central simple algebra of dimension \(2^{2l}\) (Algebra, Chap. IX, §9, no. 4, Th. 2). Let N be the exterior algebra of the maximal isotropic subspace \(F'\) generated by \(e_{-1}, \ldots, e_{-l}\) . Identify F with the dual of \(F'\) by means of \(\Psi\) and for \(x \in F'\) (resp. \(y \in F\) ) denote by \(\lambda(x)\) (resp. \(\lambda(y)\) ) the left exterior product with x (resp. the left interior product with y) in N; if \(a_{1}, \ldots, a_{k} \in F'\) , then

\[
\lambda (x). (a _ {1} \wedge \dots \wedge a _ {k}) = x \wedge a _ {1} \wedge \dots \wedge a _ {k}
\]

\[
\lambda (y). \left(a _ {1} \wedge \dots \wedge a _ {k}\right) = \sum_ {i = 1} ^ {k} (- 1) ^ {i - 1} \Psi \left(a _ {i}, y\right) a _ {1} \wedge \dots \wedge a _ {i - 1} \wedge a _ {i + 1} \wedge \dots \wedge a _ {k}.
\]

It is easily verified that \(\lambda (x)^{2} = \lambda (y)^{2} = 0\) and that

\[
\lambda (x) \lambda (y) + \lambda (y) \lambda (x) = \Psi (x, y). 1.
\]

It follows (Algebra, Chap. IX, §9, no. 1) that there exists a unique homomorphism (again denoted by \(\lambda\) ) from \(\mathrm{C}(\tilde{\mathbf{Q}})\) to \(\mathrm{End}(\mathrm{N})\) extending the map \(\lambda : \mathrm{F} \cup \mathrm{F}' \to \mathrm{End}(\mathrm{N})\) . Since \(\dim \mathrm{N} = 2^l\) and since \(\mathrm{C}(\tilde{\mathbf{Q}})\) has a unique class of simple modules, of dimension \(2^l\) (Algebra, Chap. IX, §9, no. 4, Th. 2), the representation of \(\mathrm{C}(\tilde{\mathbf{Q}})\) on \(\mathrm{N}\) defined by \(\lambda\) is irreducible and is a spinor representation of \(\mathrm{C}(\tilde{\mathbf{Q}})\) (loc. cit.).

Consider now the map \(\mu:v\mapsto e_{0}v\) from \(\tilde{\mathrm{V}}\) to \(\mathrm{C}^{+}(\mathrm{Q})\) . For \(v\in\tilde{V}\) , we have

\[
(e _ {0} v) ^ {2} = - e _ {0} ^ {2} v ^ {2} = - \mathrm{Q} (e _ {0}) \mathrm{Q} (v) = \mathrm{Q} (v) = \tilde {\mathrm{Q}} (v)
\]

and \(\mu\) extends uniquely to a homomorphism, again denoted by \(\mu\) , from \(\mathrm{C}(\tilde{\mathbb{Q}})\) to \(\mathrm{C}^{+}(\mathrm{Q})\) . Since \(\mathrm{C}(\tilde{\mathbb{Q}})\) is simple and since

\[
\dim \mathrm{C} ^ {+} (\mathrm{Q}) = \dim \mathrm{C} (\tilde {\mathrm{Q}}) = 2 ^ {2 l},
\]

we see that \(\mu\) is an isomorphism. Consequently, \(\lambda \circ \mu^{-1}\) defines a simple \(\mathrm{C}^{+}(\mathrm{Q})\) -module structure on N and \(\rho = \lambda \circ \mu^{-1} \circ f\) is an irreducible representation of g on N (Lemma 1 (vi)).

On the other hand, in view of Lemma 1 (i), we have

\[
f (H _ {i}) = \frac {1}{2} (e _ {i} e _ {- i} - e _ {- i} e _ {i}).
\]

Since \(e_{i}e_{-i} = -e_{0}^{2}e_{i}e_{-i} = e_{0}e_{i}e_{0}e_{-i}\) and \(e_{i}e_{-i} + e_{-i}e_{i} = 1\) , we have

\[
\mu^ {- 1} \circ f (H _ {i}) = \frac {1}{2} - e _ {- i} e _ {i} = - \frac {1}{2} + e _ {i} e _ {- i}.
\]

We deduce that, for \(1 \leq i_{1} < \cdots < i_{k} \leq l\) :

\[
\rho (H _ {i}) (e _ {- i _ {1}} \wedge \dots \wedge e _ {- i _ {k}}) = \left\{ \begin{array}{l l} - \frac {1}{2} e _ {- i _ {1}} \wedge \dots \wedge e _ {- i _ {k}} & \text { if } i \in \{i _ {1}, \ldots , i _ {k} \} \\ \frac {1}{2} e _ {- i _ {1}} \wedge \dots \wedge e _ {- i _ {k}} & \text { if } i \notin \{i _ {1}, \ldots , i _ {k} \} \end{array} \right.
\]

and for \(h\in \mathfrak{h}\)

\[
\begin{array}{l} \rho (h) (e _ {- i _ {1}} \wedge \dots \wedge e _ {- i _ {k}}) \\ = (\frac {1}{2} (\varepsilon_ {1} + \dots + \varepsilon_ {l}) - (\varepsilon_ {i _ {1}} + \dots + \varepsilon_ {i _ {k}})) (h) (e _ {- i _ {1}} \wedge \dots \wedge e _ {- i _ {k}}). \end{array} \tag {5}
\]

This shows that the highest weight of \(\rho\) is \(\varpi_{l}\) . We call \(\rho\) the spinor representation of g. Note that its weights are all simple (moreover, \(\varpi_{l}\) is a minuscule weight).

\begin{enumerate}
\def\labelenumi{(\Alph{enumi})}
\setcounter{enumi}{21}
\tightlist
\item
  We have \(w_0 = -1\) , so every finite dimensional simple representation of \(\mathfrak{g}\) is orthogonal or symplectic. By Chap. VI, §4, no. 5.VI, the sum of the coordinates of \(\varpi_r\) with respect to \((\alpha_1, \ldots, \alpha_l)\) is integral for \(1 \leq r \leq l - 1\) : thus, the representation \(\bigwedge^r \sigma\) is orthogonal. Moreover, it leaves invariant the extension \(\Psi_{(r)}\) of \(\Psi\) to \(\bigwedge^r V\) .
\end{enumerate}

For the spinor representation, the sum of the coordinates of \(\varpi_{l}\) with respect to \((\alpha_{1},\ldots,\alpha_{l})\) is \(\frac{1}{2}(1+\cdots+l)=\frac{l(l+1)}{4}\) (loc. cit.). Thus, it is orthogonal for \(l\equiv0\) or -1 (mod. 4) and symplectic for \(l\equiv1\) or 2 (mod. 4). In fact, consider the bilinear form \(\Phi\) on \(N=\bigwedge F'\) defined as follows: if \(x\in\bigwedge^{p}F'\) and \(y\in\bigwedge^{q}F'\) , put \(\Phi(x,y)=0\) if \(p+q\neq l\) and

\[
x \wedge y = (- 1) ^ {\frac {p (p + 1)}{2}} \varPhi (x, y) e _ {- 1} \wedge \dots \wedge e _ {- r}
\]

if \(p + q = l\) . It is easily verified that \(\Phi\) is non-degenerate and is orthogonal for \(l \equiv 0, -1 \pmod{4}\) and alternating for \(l \equiv 1, 2 \pmod{4}\) . On the other hand, in view of Lemma 1 (i),

\[
f (X _ {\varepsilon_ {i}}) = e _ {0} e _ {i}, \quad f (X _ {- \varepsilon_ {i}}) = - e _ {0} e _ {- i}
\]

for \(1 \leq i \leq l\) and

\[
f (X _ {\varepsilon_ {i} - \varepsilon_ {j}}) = \frac {1}{2} (e _ {i} e _ {- j} - e _ {- j} e _ {i}) = e _ {i} e _ {- j} = e _ {0} e _ {i} e _ {0} e _ {- j}
\]

for \(1 \leq i < j \leq l\) , and similarly

\[
f (X _ {\varepsilon_ {j} - \varepsilon_ {i}}) = - e _ {0} e _ {j} e _ {0} e _ {- i}, f (X _ {\varepsilon_ {i} + \varepsilon_ {j}}) = e _ {0} e _ {i} e _ {0} e _ {j}, f (X _ {- \varepsilon_ {i} - \varepsilon_ {j}}) = e _ {0} e _ {- i} e _ {0} e _ {- j};
\]

hence

\[
\mu^ {- 1} \circ f (X _ {\varepsilon_ {i}}) = e _ {i}, \quad \mu^ {- 1} f (X _ {- \varepsilon_ {i}}) = - e _ {- i}
\]

and

\[
\mu^ {- 1} \circ f (X _ {\pm \varepsilon_ {i} \pm \varepsilon_ {j}}) = c e _ {\pm i} e _ {\pm j} \quad \text {   for   } 1 \leq i, j \leq l, i \neq j \text {   with   } c \in \{1, - 1 \}.
\]

It is now painless to verify that \(\Phi\) is actually g-invariant (cf.~Exerc. 18).

\begin{enumerate}
\def\labelenumi{(\Roman{enumi})}
\setcounter{enumi}{5}
\tightlist
\item
  For \(x \in \mathfrak{g}\) , the characteristic polynomial of \(\sigma(x)\) takes the form
\end{enumerate}

\[
\mathrm{T} ^ {2 l + 1} + f _ {1} (x) \mathrm{T} ^ {2 l} + f _ {2} (x) \mathrm{T} ^ {2 l - 1} + \dots + f _ {2 l + 1} (x)
\]

where \(f_{1},\ldots ,f_{2l + 1}\) are invariant polynomial functions on \(\mathfrak{g}\) .

If \(x = \xi_{1}H_{1} + \cdots + \xi_{l}H_{l} \in h\) , the \(f_{i}(x)\) are, up to sign, the elementary symmetric functions of \(\xi_{1}, \ldots, \xi_{l}, -\xi_{1}, \ldots, -\xi_{l}\) ; these symmetric functions are zero in odd degrees, and

\[
\mathrm{T} ^ {2 l + 1} + f _ {2} (x) \mathrm{T} ^ {2 l - 1} + f _ {4} (x) \mathrm{T} ^ {2 l - 3} + \dots + f _ {2 l} (x) \mathrm{T} = \mathrm{T} (\mathrm{T} ^ {2} - \xi_ {1} ^ {2}) \dots (\mathrm{T} ^ {2} - \xi_ {l} ^ {2})
\]

so that \(f_{2},\ldots,f_{2l}\) are, up to sign, the elementary symmetric functions of \(\xi_{1}^{2},\ldots,\xi_{l}^{2}\) , which are algebraically independent generators of \(\mathbf{S}(\mathfrak{h}^{*})^{\mathrm{W}}\) (Chap. VI, §4, no. 5.IX). In view of §8, no. 3, Th. 1 (i), we see that \(f_{1}=f_{3}=f_{5}=\cdots=0\) and that \((f_{2},f_{4},\ldots,f_{2l})\) is an algebraically free family generating the algebra of invariant polynomial functions on g.

\begin{enumerate}
\def\labelenumi{(\Roman{enumi})}
\setcounter{enumi}{6}
\tightlist
\item
  Since the only automorphism of the Dynkin graph is the identity, we have \(\mathrm{Aut}(\mathfrak{g}) = \mathrm{Aut}_0(\mathfrak{g})\) .
\end{enumerate}

Let \(\Sigma\) be the group of similarities of V relative to \(\Psi\) . For all \(g \in \Sigma\) , let \(\varphi(g)\) be the automorphism \(x \mapsto gxg^{-1}\) of \(\mathfrak{g}\) . Then \(\varphi\) is a homomorphism from \(\Sigma\) to \(\operatorname{Aut}(\mathfrak{g})\) . We show that it is surjective. Let \(\alpha \in \operatorname{Aut}(\mathfrak{g}) = \operatorname{Aut}_0(\mathfrak{g})\) . By Prop. 2 of §7, no. 1, there exists \(s \in \mathbf{GL}(V)\) such that \(\alpha(x) = sxs^{-1}\) for all \(x \in \mathfrak{g}\) . Then \(s\) transforms \(\Psi\) into a bilinear form \(\Psi'\) on V that is invariant under \(\mathfrak{g}\) , and hence proportional to \(\Psi\) (§7, no. 5, Prop. 12). This proves that \(s \in \Sigma\) .

Since the identity representation of g is irreducible, its commutant reduces to the scalars ( \(\S6\) , no. 1, Prop. 1), so the kernel of \(\varphi\) is \(k^{*}\) . Thus, the group \(\operatorname{Aut}(\mathfrak{g}) = \operatorname{Aut}_{0}(\mathfrak{g})\) can be identified with \(\Sigma/k^{*}\) . But, it follows from Algebra, Chap. IX, \(\S6\) , no. 5, that the group \(\Sigma\) is the product of the groups \(k^{*}\) and \(\mathbf{SO}(\Psi)\) ; hence \(\operatorname{Aut}(\mathfrak{g}) = \operatorname{Aut}_{0}(\mathfrak{g})\) can be identified with \(\mathbf{SO}(\Psi)\) .

Let \(\mathbf{O}_0^+(\varPsi)\) be the reduced orthogonal group of \(\varPsi\) (Algebra, Chap. IX, §9, no. 5). Since \(\mathbf{SO}(\varPsi)/\mathbf{O}_0^+(\varPsi)\) is commutative (loc. cit.), the group \(\mathrm{Aut}_e(\mathfrak{g})\) is contained in \(\mathbf{O}_0^+(\varPsi)\) (§11, no. 2, Prop. 3); in fact, it is equal to it (Exerc. 7).

\begin{enumerate}
\def\labelenumi{(\Roman{enumi})}
\setcounter{enumi}{7}
\tightlist
\item
  The canonical bilinear form \(\Phi_{R}\) on \(h^{*}\) is given by
\end{enumerate}

\[
\varPhi_ {\mathrm{R}} (\xi_ {1} \varepsilon_ {1} + \dots + \xi_ {l} \varepsilon_ {l}, \xi_ {1} ^ {\prime} \varepsilon_ {1} + \dots + \xi_ {l} ^ {\prime} \varepsilon_ {l}) = \frac {1}{4 l - 2} (\xi_ {1} \xi_ {1} ^ {\prime} + \dots + \xi_ {l} \xi_ {l} ^ {\prime})
\]

(Chap. VI, §4, no. 5.V). The isomorphism from \(\mathfrak{h}\) to \(\mathfrak{h}^*\) defined by \(\varPhi_{\mathrm{R}}\) takes \(H_{i}\) to \((4l - 2).\varepsilon_{i}\) . Thus, the inverse form of \(\varPhi_{\mathrm{R}}\) , that is the restriction to \(\mathfrak{h}\) of the Killing form, is

\[
\varPhi (\xi_ {1} H _ {1} + \dots + \xi_ {l} H _ {l}, \xi_ {1} ^ {\prime} H _ {1} + \dots + \xi_ {l} ^ {\prime} H _ {l}) = (4 l - 2) (\xi_ {1} \xi_ {1} ^ {\prime} + \dots + \xi_ {l} \xi_ {l} ^ {\prime}).
\]

\begin{enumerate}
\def\labelenumi{(\Roman{enumi})}
\setcounter{enumi}{8}
\tightlist
\item
  Recall the \(X_{\alpha}\) ( \(\alpha \in R\) ) defined by the formulas (3). It is easy to verify that \([X_{\alpha}, X_{-\alpha}] = -H_{\alpha}\) for \(\alpha \in R\) . On the other hand, let M be the matrix \(I + E_{0,0}\) ; since \(M = S^{t}M^{-1}S\) , the map
\end{enumerate}

\[
\theta : g \mapsto - M ^ {- 1 t} g M
\]

is an automorphism of \(\mathfrak{g}\) and \(\theta(X_{\alpha}) = X_{-\alpha}\) for all \(\alpha \in \mathrm{R}\) . Consequently, \((X_{\alpha})\) is a Chevalley system in \((\mathfrak{g},\mathfrak{h})\) .

Assume that \(k = \mathbf{Q}\) . The Cartan subalgebra \(\mathfrak{h}\) has two permissible lattices: the lattice \(\mathrm{Q}(\mathrm{R}^{\vee})\) generated by the \(H_{\alpha}\) and the lattice \(\mathrm{P}(\mathrm{R}^{\vee})\) that is generated by the \(H_{i}\) and consists of the diagonal matrices in \(\mathfrak{h}\) with integer entries. It follows that \(\mathfrak{o}_S(2l + 1, \mathbf{Z})\) (the set of matrices in \(\mathfrak{g}\) with integer entries) is the Chevalley order \(\mathrm{P}(\mathrm{R}^{\vee}) + \sum \mathbf{Z}.X_{\alpha}\) in \(\mathfrak{g}\) . Since \((X_{\pm \varepsilon_i})^2 = 2E_{\pm i,\mp i}, (X_{\pm \varepsilon_i})^3 = 0\) and \((X_{\pm \varepsilon_i\pm \varepsilon_j})^2 = 0\) , we see that the lattice \(\mathcal{V}\) generated by the Witt basis \((e_i)_{-l\leq i\leq l}\) is an admissible lattice for \(\mathfrak{o}_S(2l + 1,\mathbf{Z})\) in \(V\) . The same is true for \(\bigwedge^r\mathcal{V}\) in \(\bigwedge^r V\) .

Now consider the spinor representation \(\rho\) of g on \(N = \bigwedge F'\) . As its weights do not map \(\mathrm{P}(\mathrm{R}^{\vee})\) to Z, it has no admissible lattice for \(\mathfrak{o}_{S}(2l + 1, \mathbf{Z})\) . On the other hand, the lattice N generated by the canonical basis \((e_{-i_{1}} \wedge \cdots \wedge e_{-i_{k}})\) of N (for \(1 \leq i_{1} < \cdots < i_{k} \leq l\) ) is an admissible lattice for the Chevalley order \(\mathcal{G} = \mathrm{Q}(\mathrm{R}^{\vee}) + \sum_{\alpha \in \mathrm{R}} \mathbf{Z}. X_{\alpha}\) . Indeed, it is immediate that N is stable under the exterior product with the \(e_{-i}\) and the interior product with the \(e_{i}\) (for \(1 \leq i \leq l\) ). The formulas of (V) then show that N is stable under \(\rho(\mathcal{G})\) . Moreover, since \(\rho(X_{\alpha})^{2} = 0\) for all \(\alpha \in R\) , it follows that N is admissible.

